\section{Class Diagram}

\subsection*{First-cut class diagram: entities}

Figure~\ref{fig:entities} shows the entity classes. Every entity comes from the
candidate class table in section C. User is abstract (its name and its
\texttt{role()} operation are in italics) and Customer, FlightAgent and
Administrator inherit from it. Flight is composed of its Seats, and the other
relationships are plain associations with multiplicities at both ends. Seat
availability is derived (\texttt{/ available}) from the active Reservations.
Promotion has no associations. The five enumerations hold the fixed value sets.

\diagramland{10-class-entities}{First-cut class diagram: entity classes}{fig:entities}

\subsection*{First-cut class diagram: control and boundary classes}

The diagram below adds the \texttt{<<control>>} and
\texttt{<<boundary>>} classes. Each control class comes from a use case of
section B and its operations come from the verbs of section C. Each boundary
class uses one control class, and the control classes use the entities they
coordinate. A boundary class never reaches an entity directly. The three
dashboards belong to MainFrame, which shows the promotion at start-up through
PromotionController. ConfirmationWriter is a helper that PaymentController
calls. The data access classes and the design patterns are added in the refined
diagram.

\diagram{11-class-control-boundary}{First-cut class diagram: control and boundary classes}

\paragraph{Operations of the boundary classes.} The boundary boxes are drawn
without operations to keep the diagram readable. The operations that the
sequence diagrams of section D send to a view are of two kinds: the user acts on
the view (\texttt{enter...}, \texttt{choose...}, \texttt{pick...},
\texttt{submit...}) and the controller tells the view to show something
(\texttt{show...}). They belong to the view classes below.

\begingroup
\small
\setlength{\tabcolsep}{4pt}
\begin{longtable}{@{}>{\raggedright\arraybackslash}p{4.6cm}>{\raggedright\arraybackslash}p{9.8cm}@{}}
\toprule
\textbf{Boundary class} & \textbf{Operations used in the sequence diagrams}\\
\midrule
\endhead
\ctl{FlightSearchView} & enterCriteria, showError, showMessage, showTable\\
\ctl{SeatSelectionView} & chooseBook, pickSeat, showSeatMap, showPrice, showError\\
\ctl{PaymentView} & enterDetails, showConfirmation, showError\\
\ctl{FlightManagementView} & submitFlight, chooseRemove, showFlight, showMessage, showError\\
\bottomrule
\end{longtable}
\endgroup

\subsection*{Refined class diagram}

The first-cut diagrams are refined in two steps: the business layer and the
data layer. The refinement uses inheritance (the three roles inherit from
User), realization (interfaces and the classes that implement them),
composition (a Flight is composed of its Seats) and polymorphism (the
payment methods and the reservation states), and applies six design patterns.
The attributes and associations of the entities do not change.

\paragraph{Design patterns.}

\begin{description}
\item[Singleton.] The database file must never be opened twice, so
  \texttt{DatabaseConnection} has exactly one instance, reached only through
  \texttt{getInstance()} from the composition root \texttt{AppContext}.
\item[DAO.] Controllers must not know about SQL, so every entity has a DAO
  interface in the data package, realized by one package-private SQLite class,
  and the database can be replaced without touching a business class;
  \texttt{DaoFactory} creates the DAOs.
\item[Strategy.] Credit and debit payments authorize differently, so each is a
  class behind the \texttt{PaymentMethod} interface and PaymentController
  calls \texttt{authorize} without an \texttt{if} on the type;
  \texttt{PaymentMethods.\allowbreak forType} picks the strategy.
\item[State.] A reservation reacts differently to the same call in each of its
  five states, so Reservation hands \texttt{assignSeat}, \texttt{confirmPayment},
  \texttt{changeFlight}, \texttt{cancel} and \texttt{complete} to a
  \texttt{ReservationState} object, with one class per state as in the
  reservation state diagram.
\item[Factory.] The code that logs in or registers a user must not name the
  subclasses, so \texttt{UserFactory.\allowbreak create} returns the right subclass of User
  for a role.
\item[Observer.] The monthly promotion must reach the screen without the
  controller knowing about Swing, so PromotionController notifies the
  registered \texttt{PromotionListener}s and MainFrame is the listener that
  shows the dialog.
\end{description}

\paragraph{New classes.} \texttt{AppContext} builds the controllers and the DAOs
in one place and is the only class that calls \texttt{getInstance()}.
\texttt{Authorizer} checks the role of the user and \texttt{SessionContext} holds
the logged-in user; both are used by the controllers and appear in the sequence
diagrams. The five state classes, the two payment classes and the twelve SQLite
classes exist only to realize an interface. No class outside the data package
uses a SQLite class.

Three classes are added to the data layer and the utility package.
ConfirmationWriter does file access, so the refinement moves it from the
control classes to the data package; PaymentController calls it, and so does
ReservationController when a modified reservation keeps its price. DatabaseSeeder
fills an empty database with demo data and hashes the demo passwords with
\texttt{PasswordHasher} of the utility package, which the controllers also use
for log in and registration. ChinookAir is the entry point and calls
\texttt{AppContext}.

\diagramland{12-class-refined}{Refined class diagram: business layer and design patterns}{fig:refinedbusiness}

\diagramland{13-class-data}{Refined class diagram: data layer}{fig:refineddata}

\paragraph{Control classes and their DAOs.} Controllers receive their DAO
interfaces through the constructor, so the table below is also the list of
dependencies of the data layer.

\begingroup
\small
\setlength{\tabcolsep}{4pt}
\begin{longtable}{@{}>{\raggedright\arraybackslash}p{4.6cm}>{\raggedright\arraybackslash}p{9.8cm}@{}}
\toprule
\textbf{Control class} & \textbf{DAO interfaces it uses}\\
\midrule
\endhead
\ctl{LoginController} & UserDao\\
\ctl{CustomerController} & CustomerDao, UserDao\\
\ctl{FlightSearchController} & FlightDao, SeatDao, AirportDao, AirlineDao\\
\ctl{BookingController} & FlightDao, SeatDao, ReservationDao\\
\ctl{PaymentController} & ReservationDao, PaymentDao (and ConfirmationWriter)\\
\ctl{ReservationController} & ReservationDao, PaymentDao, FlightDao, SeatDao (and ConfirmationWriter)\\
\ctl{FlightManagementController} & FlightDao, SeatDao, RouteDao, AirlineDao, AircraftDao, ReservationDao\\
\ctl{RouteController} & RouteDao, AirportDao\\
\ctl{AircraftController} & AircraftDao\\
\ctl{ScheduleController} & FlightDao\\
\ctl{PromotionController} & PromotionDao, AppStateDao\\
\bottomrule
\end{longtable}
\endgroup

\paragraph{Classes added while coding.} Writing the code added small classes that
no first-cut or refined diagram had. The first three groups are drawn in the
boxes \emph{Added while coding} of the refined diagrams; the presentation helpers
only appear in this table. None of them changes a class or an operation of the
design above. The operations that were added to existing classes are
\texttt{airports}, \texttt{airlines} and \texttt{today} on FlightSearchController
(the search screen lists airports and airlines through the controller, not a DAO),
\texttt{canCancel} and \texttt{canModify} on ReservationController (so a screen
does not repeat the rule), \texttt{isArchived} on FlightDao and
\texttt{replaceSeats} on SeatDao.

\begingroup
\small
\setlength{\tabcolsep}{4pt}
\begin{longtable}{@{}>{\raggedright\arraybackslash}p{2.9cm}>{\raggedright\arraybackslash}p{5.8cm}>{\raggedright\arraybackslash}p{5.7cm}@{}}
\toprule
\textbf{Package} & \textbf{Class} & \textbf{Why it exists}\\
\midrule
\endhead
\ctl{business.control} & \ctl{AccessDeniedException} & thrown when the role of the user may not do a task; the screen shows one plain sentence\\
\ctl{business.control} & \ctl{AuthenticationException} & thrown when a user name or password is wrong\\
\ctl{business.control} & \ctl{AircraftConflicts} & shared check of FlightManagementController and ScheduleController that an aircraft is free at a time (package-private)\\
\ctl{business.control} & \ctl{ConfirmationNumberGenerator} & interface that makes confirmation numbers, so tests can use fixed ones\\
\ctl{business.control} & \ctl{RandomConfirmationNumberGenerator} & its real implementation: CA and six random letters or digits\\
\ctl{business.control} & \ctl{Controllers} & one object with every controller and the session, handed to the screens by AppContext\\
\ctl{business.control} & \ctl{SearchCriteria} & what the search screen typed: origin, destination, date, airline\\
\ctl{business.control} & \ctl{FlightOffer} & one line of the search result: a flight, seats left, lowest price\\
\ctl{business.control} & \ctl{FlightDetails} & what the flight form typed, for add and update of a flight\\
\ctl{business.control} & \ctl{PaymentResult} & what came out of a payment: payment, reservation, confirmation file, message\\
\ctl{business.entity} & \ctl{AbstractReservationState} & base of the five state classes, so a state only writes what it allows\\
\ctl{business.entity} & \ctl{PricingPolicy} & the one place that turns a base price and a seat class into a price\\
\ctl{business.entity} & \ctl{SeatLayout} & works out the seats of an aircraft from its rows and seats per row\\
\ctl{business.entity} & \ctl{PaymentDetails} & the card details of a payment; only the masked number is kept\\
\ctl{business.entity} & \ctl{AuthorizationResult} & approved or declined, with the reason; the answer of a PaymentMethod\\
\ctl{business.entity} & \ctl{SeatUnavailableException} & thrown when another reservation holds the seat\\
\ctl{business.entity} & \ctl{StorageException} & base of the database errors, so the screens do not import the data package\\
\ctl{data} & \ctl{DataAccessException} & a database error with a message; extends StorageException\\
\ctl{data} & \ctl{ConfirmationWriteException} & the confirmation file could not be written\\
\ctl{data} & \ctl{DatabaseInitializer} & creates the tables from schema.sql (package-private)\\
\ctl{data} & \ctl{JdbcSupport} & small helpers for queries, updates and transactions (package-private)\\
\ctl{data} & \ctl{DbDates} & writes and reads dates as ISO text (package-private)\\
\ctl{presentation} & \ctl{Navigator, \allowbreak ViewId, \allowbreak Screen} & how a screen asks for another one, the list of screens, and the hook that runs when a screen is shown\\
\ctl{presentation} & \ctl{Dashboard} & base of AgentDashboard and AdminDashboard (CustomerHome is its own screen)\\
\ctl{presentation} & \ctl{BookingSession} & what is picked while booking or modifying: flight, customer, reservation, payment\\
\ctl{presentation} & \ctl{InputParser, \allowbreak ErrorMessages, \allowbreak Alert, \allowbreak ConfirmDialog} & turn typed text into values, exceptions into plain sentences, show a message, and ask the user to confirm\\
\ctl{presentation} & \ctl{SeatGridModel} & the seat map as rows and columns, without Swing\\
\ctl{presentation} & \ctl{ReservationTableModel, \allowbreak CustomerTableModel, \allowbreak PaymentTableModel, \allowbreak RouteTableModel, \allowbreak AircraftTableModel, \allowbreak ScheduleTableModel} & the rows of each table\\
\ctl{presentation} & \ctl{CustomerForm, \allowbreak PromotionPanel, \allowbreak FlightResultList} & parts of screens: the customer form (register and profile), the promotion text and the list of flights found by a search\\
\ctl{root} & \ctl{SelfTest} & the check behind \texttt{-{}-selftest}: it opens a throw-away database, logs in and searches (package-private)\\
\bottomrule
\end{longtable}
\endgroup
